\kbvchapter{ಪಾಠ ೧೧}{ನಾಮಪದಗಳ ರೂಪಬದಲಾವಣೆ ({\kbvlatin iii})}

\kbvsubtitle{ವಿಭಕ್ತಿ ಪ್ರತ್ಯಯಗಳು}

\kbvsechead{ಈ ವಾಕ್ಯಗಳನ್ನು ಓದಿ ನೋಡಿ :}

\vspace{2.5pt plus 3pt minus 0.5pt}\noindent\hspace*{1.2em}\kbvnum{(೧)}\hspace{0.55em}ರಾಮನು ದಶರಥನ ಮಗ.\par

\vspace{2.5pt plus 3pt minus 0.5pt}\noindent\hspace*{1.2em}\kbvnum{(೨)}\hspace{0.55em}ರಾಮನನ್ನು ಸೀತೆಯು ಮದುವೆಯಾದಳು.\par

\vspace{2.5pt plus 3pt minus 0.5pt}\noindent\hspace*{1.2em}\kbvnum{(೩)}\hspace{0.55em}ರಾಮನಿಂದ ಮಾರೀಚನ ಸಂಹಾರವಾಯಿತು.\par

\vspace{2.5pt plus 3pt minus 0.5pt}\noindent\hspace*{1.2em}\kbvnum{(೪)}\hspace{0.55em}ರಾಮನಿಗೆ ಬಿಲ್ಲುವಿದ್ಯೆ ಚೆನ್ನಾಗಿ ತಿಳಿದಿತ್ತು.\par

\vspace{2.5pt plus 3pt minus 0.5pt}\noindent\hspace*{1.2em}\kbvnum{(೫)}\hspace{0.55em}ರಾಮನ ದೆಸೆಯಿಂದ ಸುಗ್ರೀವನಿಗೆ ರಾಜ್ಯ ದಕ್ಕಿತು.\par

\vspace{2.5pt plus 3pt minus 0.5pt}\noindent\hspace*{1.2em}\kbvnum{(೬)}\hspace{0.55em}ರಾಮನ ಭಂಟ ಹನುಮಂತ.\par

\vspace{2.5pt plus 3pt minus 0.5pt}\noindent\hspace*{1.2em}\kbvnum{(೭)}\hspace{0.55em}ರಾಮನಲ್ಲಿ ಪ್ರಜಾವಾತ್ಸಲ್ಯ ತುಂಬಿತ್ತು.\par

\vspace{2.5pt plus 3pt minus 0.5pt}\noindent\hspace*{1.2em}\kbvnum{(೮)}\hspace{0.55em}ರಾಮನೇ !... ನೀನೇ ನಮಗೆ ದಿಕ್ಕು.\par

\indent ಮೇಲಿನ ವಾಕ್ಯಗಳಲ್ಲಿ ದಪ್ಪಕ್ಷರಗಳಲ್ಲಿರುವ ಪದಗಳನ್ನು ಗಮನಿಸಿ. ಇಲ್ಲಿ “ರಾಮ” ಎಂಬ ನಾಮಪ್ರಕೃತಿಯ ಮೇಲೆ ಪ್ರತ್ಯಯಗಳು ಸೇರಿ ನಾಮಪದಗಳಾಗಿರುವುದನ್ನು ಕಾಣಬಹುದು.\par

\begin{kbvderiv}
\vspace{4pt plus 3pt minus 1pt}
\begin{center}\begin{adjustbox}{max width=\textwidth}
{\renewcommand{\arraystretch}{1.22}\setlength{\tabcolsep}{4pt}
\begin{tabular}{l c l}
ರಾಮನು & {\color{kbvgrey}$=$} & ರಾಮ + (ಉ) \\
ರಾಮನನ್ನು & {\color{kbvgrey}$=$} & ರಾಮ + (ಅನ್ನು) \\
ರಾಮನಿಂದ & {\color{kbvgrey}$=$} & ರಾಮ + (ಇಂದ) \\
ರಾಮನಿಗೆ & {\color{kbvgrey}$=$} & ರಾಮ + (ಇಗೆ) \\
ರಾಮನ ದೆಸೆಯಿಂದ & {\color{kbvgrey}$=$} & ರಾಮನ + (ದೆಸೆಯಿಂದ) \\
ರಾಮನ & {\color{kbvgrey}$=$} & ರಾಮ + (ಅ) \\
ರಾಮನಲ್ಲಿ & {\color{kbvgrey}$=$} & ರಾಮ + (ಅಲ್ಲಿ) \\
ರಾಮನೇ & {\color{kbvgrey}$=$} & ರಾಮ + (ಏ) \\
\end{tabular}}
\end{adjustbox}\end{center}
\vspace{4pt plus 3pt minus 1pt}
\end{kbvderiv}

\indent ಆವರಣಗಳಲ್ಲಿರುವ ಪ್ರತ್ಯಯಗಳು ಕ್ರಮವಾಗಿ ಮೇಲಿನ ವಾಕ್ಯಗಳಲ್ಲಿ ಕರ್ತೃ, ಕರ್ಮ, ಕ್ರಿಯಾಪದಗಳಿಗಿರುವ ಸಂಬಂಧವನ್ನು ಅರ್ಥವತ್ತಾಗಿ ತಿಳಿಸುತ್ತವೆ. ಈ ರೀತಿಯಾಗಿ, ನಾಮಪ್ರಕೃತಿಯ ಮುಂದೆ ಸೇರುವ ಪ್ರತ್ಯಯಗಳಿಗೆ “ವಿಭಕ್ತಿ ಪ್ರತ್ಯಯಗಳೆಂದು ಹೆಸರು.\par
